\documentclass[12pt,a4paper]{article}
\usepackage{natbib}
\usepackage{amsmath}
\usepackage[T1]{fontenc}
\usepackage[utf8]{inputenc}

\title{My First LaTeX Document}
\author{Van Meet}
\date{\today}

\begin{document}

\maketitle

\tableofcontents

\section{Hello World}

Hello World!

This is my first LaTeX document. The document is created using the
\texttt{article} class and compiled to produce a formatted output.

\section{My First LaTeX Document}

LaTeX documents are written using commands and environments. A section
is created with the \texttt{\textbackslash section} command, while a
subsection is created with \texttt{\textbackslash subsection}.

\subsection{Numbered List}

The following is an example of a numbered list:

\begin{enumerate}
    \item Create a new document.
    \item Write the document content in the editor.
    \item Save the file with the \texttt{.tex} extension.
    \item Compile the document and view the output.
\end{enumerate}

\subsection{Document Options}

The document class options set the font size and paper size. This
document uses a 12-point font and A4 paper:

\[
\texttt{\textbackslash documentclass[12pt,a4paper]\{article\}}
\]

The \texttt{\textbackslash tableofcontents} command generates the
table of contents from the document's section headings.

\section{Science!}

Einstein's famous mass-energy equivalence equation is
\(E = mc^2\). The relationship between mass and energy was introduced
by Einstein \citep{einstein1905}.

\section{Equations and Citations}

The time-dependent Schr\"odinger equation can be written as

\begin{equation}
i\hbar\frac{\partial\Psi(\mathbf{r},t)}{\partial t}
=
\left[
-\frac{\hbar^2}{2\mu}\nabla^2
+V(\mathbf{r},t)
\right]\Psi(\mathbf{r},t).
\label{eq:schrodinger}
\end{equation}

Equation~\ref{eq:schrodinger} is the Schr\"odinger equation.
Schr\"odinger's 1926 work developed wave mechanics
\citep{schrodinger1926}.

Einstein's paper is cited using the citation key
\texttt{einstein1905}, and Schr\"odinger's paper is cited using
\texttt{schrodinger1926}. The entries for both papers are stored in
the accompanying \texttt{physics.bib} file.

\bibliographystyle{agu}
\bibliography{physics}

\end{document}
